%!TEX root = main.tex
%=========================================================
\newpage

\section{Global broadcast in heterogeneous MANET (draft of proof)}

Previous works in broadcast protocols assume that every node in a network follows a unique protocol.
To the best of our knowledge, emergent overlays is the first adaptive approach that combines two state-of-the-art broadcast protocols.
This combination result in having an heterogeneous MANET where nodes switch from CF and MPR (or vice-versa).
We need to guarantee global broadcast in such heterogeneous networks.
This section contains a first step towards a general poof of full coverage in heterogeneous MANETs where nodes switch from CF and MPR to broadcast.

\textbf{Lemma 1}. Let N be a MANET where nodes have no motion and they chose in an arbitrarily way between CF or MPR as broadcast protocol. We then state that every node in N receive broadcast messages from every source node within such network.


Before giving a proof of Lemma 1, we rely on the following assumptions:

\begin{itemize}
	\item The network N is deployed assuming ideal conditions of communication, in other words, every node located at the broadcast area of relays receive broadcast messages successfully.
	\item Individual protocols guarantee full coverage. Global coverage is guaranteed in a network where every node use a unique protocol to broadcast.
	\item Border nodes have been discovered with the interoperable mechanism presented in section~\ref{sec:mechanisms}. Recall that border nodes are simply bridges between CF and MPR (or \emph{vice-versa}).
\end{itemize}

We proceed with the proof of Lemma 1 by showing that there is no node in the static network N that misses broadcast messages. Observe that this is similar to say that global broadcast takes place in N.
Given that nodes select arbitrarily the protocol to broadcast, we separate such selection in two possible outcomes.

First, assume that all nodes chose the same protocol P to broadcast.
As part of our assumptions, P guarantees global broadcast in an homogeneous network (where nodes run a unique protocol), which makes Lemma 1 a valid statement.

In the second outcome, some nodes in N have arbitrarily chosen CF as broadcast protocol and others MPR.
Let \emph{m} be a broadcast message.
Assume that there exist a node \emph{n} in range of another node \emph{n}' such that \emph{n} does not receive a copy of \emph{m}, even if \emph{n'} holds a copy of such message.
Consider now that the broadcast protocol running at \emph{n} differs from the one running at \emph{n'}, once again, we have two possible outcomes:

\begin{itemize}
	\item \textbf{\emph{n'} runs MPR and \emph{n} runs CF.} Node \emph{n'} maintains an overlay, when \emph{n'} acts as a relay this node forwards \emph{m} and eventually, \emph{n} receives a copy of such message. In case \emph{n'} is a pure receiver at the overlay, its neighbor \emph{n} have bootstrapped the interoperability mechanism (see section~\ref{sec:mechanisms}) and have chosen \emph{n'} as border node, because \emph{n} is at the edge of a system running CF. Border nodes are relays, meaning that \emph{n} will eventually receive \emph{m}.
	\item \textbf{\emph{n'} runs CF and \emph{n} runs MPR.}
\end{itemize}

We conclude that assuming the existence of a node \emph{n} that do not hear from message \emph{m}, is a false statement, which makes Lemma 1 a valid statement.


%Assume that there exist a node n within range of another node n' that do not receive a broadcast message even if n' already have a copy of such message.

